% Central Paper C content (journal-neutral). Edit here; journal wrappers in ../journals/ only \input these files.
% Every number traces to ../../FINDINGS.md.
\section{Introduction}\label{sec:intro}

The Journal Impact Factor (JIF) counts the citations a journal receives in one year to the items it published in the
two previous years. Because the JIF still shapes hiring, funding and journal choice, some journals inflate it, either
by citing themselves or by receiving concentrated citations from a partner journal, a practice known as citation
stacking \citep{heneberg2016stacking,vannoorden2013brazil,wilhite2012coercive,chorus2016coercive}. Each year Clarivate
withholds the JIF of journals whose citation patterns it judges anomalous and publishes the list, with key figures
such as the share of self-citations in the JIF numerator or the share of a recipient's numerator that comes from one
donor \citep{moussa2022suppressed,taskin2021selfcitation}.

These decisions rest on Web of Science data, which outsiders cannot inspect. Open citation data now make independent
checks possible in principle: OpenCitations releases citation links under open licences
\citep{peroni2020opencitations,heibi2019coci}, and its Meta collection gives the bibliographic records behind them
\citep{massari2024meta}. Whether these data are good enough for integrity work is less clear. Coverage comparisons
show that open sources approach the large commercial indexes on average
\citep{martinmartin2021coverage,visser2021comparison,culbert2025openalex}, but an average says little about the small,
specific sets of journals that suppressions concern.

The best-known open method for this task is CIDRE \citep{kojaku2021cidre}. It removes citations that a degree-corrected
stochastic block model can explain, scores journals as donors and recipients of the remaining ``excessive''
citations, and keeps groups of journals that pass a threshold. On Microsoft Academic Graph data it found 12 of 22
groups that Clarivate had suspended, and community detection found none.

Three questions are addressed. First, how faithfully do fully open data reproduce the figures Clarivate publishes, and where
and why do they fail? Second, does CIDRE, run as published, detect the stacking cases Clarivate listed after its
own validation period, and if not, why? Third, what do the excess citations look like at the level of individual
papers? The contributions are:
\begin{itemize}
  \item a primary-source label set of Clarivate suppressions for data years 2015--2023 and an audit of open-data
        fidelity against Clarivate's published self-citation and stacking figures, with the failures traced to
        publishers' reference deposit and to a data-citation artifact;
  \item a diagnosis of a structural blind spot in CIDRE: its group-building step drops large donor journals, so the
        pattern Clarivate most often reports, a large journal feeding a small one, cannot form a group;
  \item a recipient-centric reading of CIDRE's own excessive-edge test that flags listed stacking journals well above
        chance, also against size-matched controls and plain baselines, sometimes years before the listing, with a
        pre-specified evidence check of its flags and a prospective check on the 2024--2025 lists;
  \item an analysis of the sensitivity of CIDRE to its stochastic null-model partition, and a paper-level anatomy of
        flagged journal pairs.
\end{itemize}
All analyses run on public data; code, run logs and checksums are released (Section~\ref{sec:availability}).

